\question{Is secure aggregation reachable, and what exactly does it guarantee?}
\label{sec:secagg}
\statusMethod

\begin{shortanswer}
It is reachable, and it is now implemented and verified end-to-end: with
\texttt{secure-aggregation} enabled, a full round of training and evaluation completes across twelve
practices and the server obtains only the sum. Three limits must travel with any claim made about
it. The guarantee is \textbf{semi-honest} --- it protects against a server that follows the protocol
and inspects what it receives, not against one that deviates. \textbf{Participation stays visible}:
the server always learns who took part. And \textbf{the aggregate is still model parameters}, so
secure aggregation does not by itself make the released model anonymous.
\end{shortanswer}

\subsection{The mechanism}

Secure aggregation \cite{bonawitz2017secagg} has each pair of participating practices derive a
shared secret mask. Practice $i$ adds the mask to its update and practice $j$ subtracts it, so the
masks cancel exactly when the contributions are summed --- and only then. The server sees each
practice's masked vector, which on its own is indistinguishable from noise, and recovers a
meaningful value only from the total.

Because a practice that drops out mid-round would leave its masks uncancelled, each practice also
distributes Shamir secret shares of its own seed. If enough surviving practices contribute shares,
the dropped practice's mask can be reconstructed and removed. Two parameters govern this and are
exposed as configuration: the number of shares issued per practice, and the reconstruction threshold
--- the number required to recover a mask. The threshold is the security parameter: it is the number
of practices that would have to collude \emph{with} the server to unmask an individual contribution.

\subsection{Why it is a second code path}

\begin{table}[htbp]
  \centering
  \caption{Secure-aggregation availability, probed against the installed package rather than
  assumed from documentation.}
  \begin{tabular}{lp{5.2cm}}
    \toprule
    Probe & Result \\
    \midrule
    \texttt{secaggplus\_mod} in \texttt{flwr.clientapp.mod} (Message API) & absent \\
    \texttt{secaggplus\_mod} in \texttt{flwr.client.mod} (legacy) & present \\
    \texttt{SecAggPlusWorkflow} importable & present \\
    \texttt{SecAggPlusWorkflow} requires \texttt{LegacyContext} & yes \\
    DP mods in \texttt{flwr.clientapp.mod} & \texttt{fixedclipping\_mod}, \texttt{adaptiveclipping\_mod}, \texttt{LocalDpMod} \\
    Implemented and verified in this repository & yes \\
    \bottomrule
\end{tabular}

  \label{tab:secagg}
\end{table}

In Flower \FlwrVersion{}, secure aggregation ships only in the \emph{legacy} namespaces: the client
modifier is absent from the current message-API namespace, and the workflow that drives it demands a
legacy context. It therefore cannot compose with the strategy entry point this project's server
otherwise uses. It \emph{is} reachable by driving the legacy workflow from inside the modern server
application, and that is what is now built.

Two implementation notes, recorded because they are the kind of detail that is expensive to
rediscover:

\begin{itemize}
  \item \textbf{One client application serves both paths.} The legacy workflow speaks a different
        record shape than the message API, so the practice's handlers accept either and translate
        through Flower's own compatibility bridge. There is no duplicated training code.
  \item \textbf{The switch is ordinary run configuration, not an environment variable.} Client
        modifiers are attached at construction time, before any run configuration exists, so secure
        aggregation and local DP are each installed as a small dispatching modifier that reads the
        configuration when it is called. This also makes them survive the process boundary in
        simulation, where client applications run in worker processes that do not inherit the
        launching environment.
\end{itemize}

\subsection{The three limits, stated plainly}

\begin{enumerate}
  \item \textbf{The threat model is semi-honest.} The guarantee holds against a server that follows
        the protocol but inspects what it receives. It is \emph{not} a guarantee against a server
        that actively deviates --- most obviously by running a round with a single practice, whose
        ``aggregate'' is that practice's update. The available control is the minimum participant
        count, which is set to the full practice count so that such a round fails rather than
        quietly succeeding. That is a configuration discipline, not a cryptographic guarantee, and
        counsel should treat it as such.
  \item \textbf{Participation stays visible.} The server always learns which practices took part in
        each round, and the aggregate. Only the individual contribution is hidden. Where
        participation itself is sensitive --- a practice's presence in a CKD study may be --- secure
        aggregation does not address it.
  \item \textbf{The aggregate is still model parameters.} Secure aggregation answers \emph{who may
        see one practice's update}. It says nothing about what can be inferred from the model that
        is eventually released. That is the province of \S\ref{sec:budget} and \S\ref{sec:anontest}.
\end{enumerate}

\subsection{One protocol it cannot protect}

Secure aggregation masks numeric vectors by arithmetic that cancels on summation. The tree-based
protocol transmits serialized decision trees, which are not a numeric vector and cannot be summed;
there is nothing to mask. \textbf{Any secure-aggregation deployment is necessarily a
logistic-regression deployment.} As \S\ref{sec:modelchoice} shows, the accuracy evidence points the
same way independently, so this constraint costs nothing.
